\documentclass[12pt,letterpaper]{article}
\usepackage{graphicx,textcomp}
\usepackage{natbib}
\usepackage{setspace}
\usepackage{fullpage}
\usepackage{color}
\usepackage[reqno]{amsmath}
\usepackage{amsthm}
\usepackage{fancyvrb}
\usepackage{amssymb,enumerate}
\usepackage[all]{xy}
\usepackage{endnotes}
\usepackage{lscape}
\newtheorem{com}{Comment}
\usepackage{float}
\usepackage{hyperref}
\newtheorem{lem} {Lemma}
\newtheorem{prop}{Proposition}
\newtheorem{thm}{Theorem}
\newtheorem{defn}{Definition}
\newtheorem{cor}{Corollary}
\newtheorem{obs}{Observation}
\usepackage[compact]{titlesec}
\usepackage{dcolumn}
\usepackage{tikz}
\usetikzlibrary{arrows}
\usepackage{multirow}
\usepackage{xcolor}
\newcolumntype{.}{D{.}{.}{-1}}
\newcolumntype{d}[1]{D{.}{.}{#1}}
\definecolor{light-gray}{gray}{0.65}
\usepackage{url}
\usepackage{listings}
\usepackage{color}

\definecolor{codegreen}{rgb}{0,0.6,0}
\definecolor{codegray}{rgb}{0.5,0.5,0.5}
\definecolor{codepurple}{rgb}{0.58,0,0.82}
\definecolor{backcolour}{rgb}{0.95,0.95,0.92}

\lstdefinestyle{mystyle}{
	backgroundcolor=\color{backcolour},   
	commentstyle=\color{codegreen},
	keywordstyle=\color{magenta},
	numberstyle=\tiny\color{codegray},
	stringstyle=\color{codepurple},
	basicstyle=\footnotesize,
	breakatwhitespace=false,         
	breaklines=true,                 
	captionpos=b,                    
	keepspaces=true,                 
	numbers=left,                    
	numbersep=5pt,                  
	showspaces=false,                
	showstringspaces=false,
	showtabs=false,                  
	tabsize=2
}
\lstset{style=mystyle}
\newcommand{\Sref}[1]{Section~\ref{#1}}
\newtheorem{hyp}{Hypothesis}


\title{Problem Set 4}
\date{Due: November 27, 2025}
\author{Applied Stats/Quant Methods 1}


\begin{document}
	\maketitle
	\section*{Instructions}
	\begin{itemize}
		\item Please show your work! You may lose points by simply writing in the answer. If the problem requires you to execute commands in \texttt{R}, please include the code you used to get your answers. Please also include the \texttt{.R} file that contains your code. If you are not sure if work needs to be shown for a particular problem, please ask.
		\item Your homework should be submitted electronically on GitHub.
		\item This problem set is due before 23:59 on Thursday November 27, 2025. No late assignments will be accepted.
	\end{itemize}



	\vspace{.5cm}
\section*{Question 1: Economics}
\vspace{.25cm}
\noindent 	
In this question, use the \texttt{prestige} dataset in the \texttt{car} library. First, run the following commands:

\begin{verbatim}
install.packages(car)
library(car)
data(Prestige)
help(Prestige)
\end{verbatim} 


\noindent We would like to study whether individuals with higher levels of income have more prestigious jobs. Moreover, we would like to study whether professionals have more prestigious jobs than blue and white collar workers.

\newpage
\begin{enumerate}
	
	\item [(a)]
	Create a new variable \texttt{professional} by recoding the variable \texttt{type} so that professionals are coded as $1$, and blue and white collar workers are coded as $0$ (Hint: \texttt{ifelse}).
			\lstinputlisting[language=R, firstline=39, lastline=43]{PS04_RS.R}  
	
	\vspace{0.5cm}
	
	
	\item [(b)]
	Run a linear model with \texttt{prestige} as an outcome and \texttt{income}, \texttt{professional}, and the interaction of the two as predictors (Note: this is a continuous $\times$ dummy interaction.)
	\begin{table}[!htbp] \centering 
		\caption{Linear Model of Prestige on Income, Professional Status, and Interaction} 
		\label{tab:prestige_model} 
		\begin{tabular}{@{\extracolsep{5pt}}lc} 
			\\[-1.8ex]\hline 
			\hline \\[-1.8ex] 
			& \multicolumn{1}{c}{\textit{Dependent variable:}} \\ 
			\cline{2-2} 
			\\[-1.8ex] & Prestige \\ 
			\hline \\[-1.8ex] 
			Income & 0.003$^{***}$ \\ 
			& (0.0005) \\ 
			Professional & 37.781$^{***}$ \\ 
			& (4.248) \\ 
			Income × Professional & $-$0.002$^{***}$ \\ 
			& (0.001) \\ 
			Constant & 21.142$^{***}$ \\ 
			& (2.804) \\ 
			\hline \\[-1.8ex] 
			Observations & 98 \\ 
			R$^{2}$ & 0.787 \\ 
			Adjusted R$^{2}$ & 0.780 \\ 
			Residual Std. Error & 8.012 (df = 94) \\ 
			F Statistic & 115.878$^{***}$ (df = 3; 94) \\ 
			\hline 
			\hline \\[-1.8ex] 
			\textit{Note:}  & \multicolumn{1}{r}{$^{*}$p$<$0.1; $^{**}$p$<$0.05; $^{***}$p$<$0.01} \\ 
		\end{tabular} 
	\end{table} 
	\lstinputlisting[language=R, firstline=45, lastline=58]{PS04_RS.R}  
	
	\vspace{0.5cm}
	\item [(c)]
	Write the prediction equation based on the result.
	\lstinputlisting[language=R, firstline=61, lastline=62]{PS04_RS.R}  
	\textbf{\[
		\widehat{\text{prestige}}
		= 21.1
		+ 0.0031 \cdot \text{income}
		+ 37.78 \cdot \text{professional}
		- 0.0023 \cdot (\text{income} \times \text{professional})
		\]
	}
	
\newpage
	\item [(d)]
	Interpret the coefficient for \texttt{income}.
	
	\textbf{For non-professional occupations, each additional 1 dollar increase in average income is associated with an estimated 0.0032-point increase in occupational prestige, holding professional status constant.}
	
	\vspace{0.5cm}	
	\item [(e)]
	Interpret the coefficient for \texttt{professional}.
	
	\textbf{Holding income at 0, professional occupations are estimated to have about 37.78 points higher prestige than non-professional occupations.}
	
	\newpage
	\item [(f)]
	What is the effect of a \$1,000 increase in income on prestige score for professional occupations? In other words, we are interested in the marginal effect of income when the variable \texttt{professional} takes the value of $1$. Calculate the change in $\hat{y}$ associated with a \$1,000 increase in income based on your answer for (c).
	
	\textbf{Income coefficient + interaction term = 0.0008452. This is the increase in prestige per dollar increase in income. Therefore, with a thousand dollar increase the change in $\hat{y}$ = 0.8452 }
	
	
	\vspace{0.5cm}
	
	
	\item [(g)]
	What is the effect of changing one's occupations from non-professional to professional when her income is \$6,000? We are interested in the marginal effect of professional jobs when the variable \texttt{income} takes the value of $6,000$. Calculate the change in $\hat{y}$ based on your answer for (c).
	
	\textbf{Using the formula for the marginal effect of a dummy variable in an interaction model: \[
		\text{ME}_{\text{professional}}
		= \beta_{2} + \beta_{3} \cdot \text{income}
		\]
	and plugging in the values: 
	\[
	\text{ME}_{\text{professional}}
	= 37.7813 - 0.0023257 \times 6000
	= 23.8271
	\]
	We get 23.83, therefore, Switching from a non-professional to a professional occupation at an income of 6,000 dollars is associated with an estimated increase of about 23.83 points in occupational prestige.}
	
	
\end{enumerate}

\newpage

\section*{Question 2: Political Science}
\vspace{.25cm}
\noindent 	Researchers are interested in learning the effect of all of those yard signs on voting preferences.\footnote{Donald P. Green, Jonathan	S. Krasno, Alexander Coppock, Benjamin D. Farrer,	Brandon Lenoir, Joshua N. Zingher. 2016. ``The effects of lawn signs on vote outcomes: Results from four randomized field experiments.'' Electoral Studies 41: 143-150. } Working with a campaign in Fairfax County, Virginia, 131 precincts were randomly divided into a treatment and control group. In 30 precincts, signs were posted around the precinct that read, ``For Sale: Terry McAuliffe. Don't Sellout Virgina on November 5.'' \\

Below is the result of a regression with two variables and a constant.  The dependent variable is the proportion of the vote that went to McAuliff's opponent Ken Cuccinelli. The first variable indicates whether a precinct was randomly assigned to have the sign against McAuliffe posted. The second variable indicates
a precinct that was adjacent to a precinct in the treatment group (since people in those precincts might be exposed to the signs).  \\

\vspace{.5cm}
\begin{table}[!htbp]
	\centering 
	\textbf{Impact of lawn signs on vote share}\\
	\begin{tabular}{@{\extracolsep{5pt}}lccc} 
		\\[-1.8ex] 
		\hline \\[-1.8ex]
		Precinct assigned lawn signs  (n=30)  & 0.042\\
		& (0.016) \\
		Precinct adjacent to lawn signs (n=76) & 0.042 \\
		&  (0.013) \\
		Constant  & 0.302\\
		& (0.011)
		\\
		\hline \\
	\end{tabular}\\
	\footnotesize{\textit{Notes:} $R^2$=0.094, N=131}
\end{table}

\vspace{.5cm}
\begin{enumerate}
	\item [(a)] Use the results from a linear regression to determine whether having these yard signs in a precinct affects vote share (e.g., conduct a hypothesis test with $\alpha = .05$).
	
	\[
	\textbf{Hypothesis Test for the Effect of Yard Signs (}\alpha = 0.05\textbf{)}
	\]
	
	\textbf{Given:}
	\[
	\hat{\beta} = 0.042, \qquad SE = 0.016
	\]
	
	\textbf{1. Hypotheses}
	\[
	H_{0}:\ \beta = 0 \qquad\text{(no effect)}
	\]
	\[
	H_{1}:\ \beta \neq 0 \qquad\text{(an effect)}
	\]
	
	\textbf{2. Test Statistic}
	\[
	t = \frac{\hat{\beta}}{SE}
	= \frac{0.042}{0.016}
	= 2.625.
	\]
	
	Degrees of freedom:
	\[
	df = N - k = 131 - 3 = 128,
	\]
	where \(k = 3\) represents two predictors and an intercept.
	
	\textbf{3. p-value (two-tailed)}
	\[
	p \approx 0.0087.
	\]
	
	\textbf{4. 95\% Confidence Interval}
	\[
	0.042 \pm 1.98(0.016)
	= (0.0103,\ 0.0737).
	\]
	
	\textbf{5. Conclusion}
	
	Since
	\[
	t = 2.625 > 1.98 \quad\text{and}\quad p < 0.05,
	\]
	we reject the null hypothesis. There is a statistically significant effect of being assigned yard signs. Precincts where the anti-McAuliffe signs were posted experienced, on average, a \textbf{4.2 percentage-point increase} in the vote share for Cuccinelli, with a 95\% confidence interval of approximately 1.03 to 7.37 percentage points.
	
	
	\newpage		
	\item [(b)]  Use the results to determine whether being
	next to precincts with these yard signs affects vote
	share (e.g., conduct a hypothesis test with $\alpha = .05$).
	
	\[
	\textbf{Hypothesis Test for the Effect of Being Adjacent to Yard Signs (}\alpha = 0.05\textbf{)}
	\]
	
	\textbf{Given:}
	\[
	\hat{\beta}_{\text{adjacent}} = 0.042, \qquad SE = 0.013
	\]
	
	\textbf{1. Hypotheses}
	\[
	H_{0}:\ \beta_{\text{adjacent}} = 0 \qquad\text{(no effect)}
	\]
	\[
	H_{1}:\ \beta_{\text{adjacent}} \neq 0 \qquad\text{(an effect)}
	\]
	
	\textbf{2. Test Statistic}
	\[
	t = \frac{\hat{\beta}_{\text{adjacent}}}{SE}
	= \frac{0.042}{0.013}
	\approx 3.231.
	\]
	
	Degrees of freedom:
	\[
	df = N - k = 131 - 3 = 128,
	\]
	where \(k=3\) (two predictors plus an intercept).
	
	\textbf{3. p-value (two-tailed)}
	\[
	p \approx 0.00157.
	\]
	
	\textbf{4. 95\% Confidence Interval}
	Using the critical \(t\)-value \(t_{0.975,128}\approx 1.9787\),
	\[
	0.042 \pm 1.9787(0.013)
	= (0.01628,\ 0.06772).
	\]
	
	\textbf{5. Conclusion}
	
	Since
	\[
	t \approx 3.231 > 1.98 \quad\text{and}\quad p \approx 0.00157 < 0.05,
	\]
	we reject the null hypothesis. There is a statistically significant effect of being adjacent to precincts with the anti-McAuliffe yard signs: precincts adjacent to treatment precincts experienced, on average, an estimated \textbf{0.042 increase} in the vote share for Cuccinelli, with a 95\% confidence interval of approximately \textbf{0.0163 to 0.0677}.
	
	
	\vspace{7cm}
	\item [(c)] Interpret the coefficient for the constant term substantively.
	\vspace{0.1cm}
	
	\textbf{The constant term in the regression is 0.302. Therefore, In precincts that were not assigned yard signs and were not adjacent to any precinct with yard signs, the expected proportion of the vote going to Cuccinelli is 0.302 (i.e., 30.2 percent).}
	
	\item [(d)] Evaluate the model fit for this regression.  What does this	tell us about the importance of yard signs versus other factors that are not modeled?
	
 
 The regression reports an 
  \[
  R^{2} = 0.094,
  \]
  which means that the model explains only about 9.4\% of the variation in vote share across precincts.
  
  This low level of explained variance indicates that, although the coefficients on the treatment and adjacency variables are statistically significant, these variables account for only a small portion of the overall differences in voting outcomes. Most of the variation in vote share is driven by other factors that are not included in the model, such as demographic characteristics, partisanship, historical voting patterns, and other campaign activities.
  
  
	
\end{enumerate}  


\end{document}
